\section{Classical Virtualization Theory}
According to the papers written by Popek and Goldberg
back in 1974 \cite{citation_17}, system software must meet three
requirements in order to qualify as a VMM:
\begin{enumerate}
    \item \textbf{Faithfulness}. Except for timing effects, software
    on the VMM runs just like it does on hardware.
    \item \textbf{Performance}. Without the VMM's assistance,
    the hardware carries out the vast majority of
    guest commands.
    \item \textbf{Security}. All hardware resources are managed
    by the VMM. Trap-and-emulate, a specific
    VMM implementation approach, was so
    common in 1974 that it was thought to be the
    only workable virtualisation technique.
  \end{enumerate}
Over the years, there has been considerable confusion
caused by the informal linking of what can be virtualized
with the capacity to utilise trap-and-emulate, even though
Popek and Goldberg did not rule out the use of other
approaches. To avoid this misunderstanding, an
architecture that can be virtualised using only trap-andemulate will be referred to as classically virtualizable.
According to Popek and Goldberg's criteria, x86 can be
virtualised using the methods outlined in Section 2, even
though it is not traditionally virtualizable in this sense.
The key concepts from traditional VMM
implementations, de-privileging, shadow structures, and
traces, are reviewed in this section. 
\subsection{De-privileging}
It is possible to make any instructions that read or write
privileged state to trap when they are run in an
unprivileged context in a classically virtualizable
architecture. The VMM protects structures that the
instructions reach, such as the address range of a
memory-mapped I/O device, might occasionally cause the
traps, as can the instruction type itself, such as an out
instruction. Guest operating systems are directly executed
by a traditional VMM, although with fewer privileges.
The VMM emulates the trapping instruction against the
virtual machine state after intercepting traps from the deprivileged guest. It is simple to confirm that the final
VMM satisfies the Popek and Goldberg requirements,
and this method has been thoroughly documented in the
literature \cite{citation_10}\cite{citation_18}\cite{citation_19_book}. 
\subsection{Primary and shadow structures}
A virtual system's privileged state is by definition
different from the underlying hardware. Despite this
difference, the VMM's primary role is to deliver an
execution environment that satisfies the guest's
expectations. The VMM uses guest-level primary
structures to derive shadow structures in order to achieve
this. The VMM keeps a picture of the guest register and
refers to that image in instruction emulation as the guest
operations trap. This makes handling on CPU privileged
state, such as the page table pointer register or processor
status register, simple. On the other hand, page tables and
other off CPU privileged data might be stored in RAM. In
this situation, trapping instructions might not always
correspond with guest accesses to the privileged state. \newline \newline
For instance, because they encode mappings and
permissions, guest page table entries (gPTEs) are in a
privileged state. Every guest virtual memory reference is
dependent on the permissions and mappings stored in the
relevant gPTE, so dependencies on this privileged state
are free of traps. Any store in the guest instruction stream
has the ability to change this in memory privileged state,
or it may even be changed implicitly as a result of a DMA
I/O transaction. Similar challenges arise with memory
mapped I/O devices, since reads and writes to this
privileged data might come from nearly any memory
operation in the guest instruction stream. 
\subsection{Memory traces}
VMMs usually use hardware page protection methods to
trap access to in memory primary structures in order to
preserve coherency of shadow structures. For example,
guest PTEs may be write protected if shadow PTEs have
been created for them. In general, memory mapped
devices need to be protected for both writing and reading.
This technology of page protection is called tracing. Like
privileged instruction faults, trace faults are handled by
classical VMMs by decoding the guest instruction that is
causing the fault, simulating its impact in the primary
structure, and then propagating the change to the shadow
structure.
\subsection{Hardware VMM implementation}
VMware’s hardware VMM's (or any VMM adopting the
extensions) structure is fundamentally prescribed by the
hardware extensions, which offer a comprehensive
virtualisation solution. When executing a guest in
protected mode, the VMM runs vmrun after filling in a
VMCB with the guest's current state. When a guest quits,
the VMM reads the VMCB fields that specify the exit
conditions and vectors to the relevant emulation code. A
significant portion of this emulation code is shared with
the VMM program. Peripheral device models, guest
interrupt delivery code, and several infrastructure
functions like logging, synchronisation, and host OS
interaction are all included. The hardware VMM also
inherits the software VMM's use of the shadowing
mechanism outlined in Section 5.2, as current
virtualisation hardware does not explicitly enable MMU
virtualisation. 


